\begin{abstract}
% 是否需要加卵巢癌约束，摘要首句以及标题
In the treatment of ovarian cancer, precise residual disease prediction is significant for clinical and surgical decision-making.
However, traditional methods are either invasive~(e.g., laparoscopy) or time-consuming~(e.g., manual analysis).
Recently, deep learning methods make many efforts in automatic analysis of medical images.
%Recently, deep learning techniques have made significant improvement in the automatic analysis of medical images.
Despite the remarkable progress, most of them underestimated the importance of 3D image information of disease, which might brings a limited performance for residual disease prediction, especially in small-scale datasets.
%Despite the remarkable progress, the majority of them underestimated the importance of disease-related 3D image information, which might brings a limited performance for residual disease prediction, especially in small-scale datasets.
To this end, in this paper,
we propose a novel Multi-View Attention Learning~(MuVAL) method for residual disease prediction,
which focuses on the comprehensive learning of 3D Computed Tomography (CT) images in a multi-view manner.
Specifically, we first obtain multi-view of 3D CT images from transverse, coronal and sagittal views.
To better represent the image features in a multi-view manner, we further leverage attention mechanism to help find the more relevant slices in each view. 
Extensive experiments on a dataset of 111 patients show that our method outperforms existing deep-learning methods. 

\end{abstract}